\documentclass[11pt]{amsart}
\usepackage[T1]{fontenc}
\usepackage{lmodern}
\usepackage[margin=1in]{geometry}
\usepackage{amsmath,amssymb,amsthm,mathtools}
\usepackage{booktabs}
\usepackage[hidelinks]{hyperref}
\newtheorem{theorem}{Theorem}[section]
\newtheorem{corollary}[theorem]{Corollary}
\newtheorem{proposition}[theorem]{Proposition}
\newtheorem{lemma}[theorem]{Lemma}
\theoremstyle{definition}
\newtheorem{definition}[theorem]{Definition}
\newtheorem{remark}[theorem]{Remark}
\newcommand{\kk}{k}
\newcommand{\CI}{\operatorname{CI}}
\newcommand{\coker}{\operatorname{coker}}
\newcommand{\Hilb}{\operatorname{Hilb}}
\newcommand{\ch}{\operatorname{char}}
\newcommand{\R}{\mathbb R}
\newcommand{\Q}{\mathbb Q}
\DeclareMathOperator{\rank}{rank}
\title[A strong Lefschetz family]{A strong Lefschetz family among equigenerated monomial almost complete intersections}
\author{Ke-kuk Chow}
\date{7 November 2026}
\subjclass[2020]{Primary 13E10; Secondary 13C40, 13F55}
\keywords{Strong Lefschetz property, monomial almost complete intersection, graded Jordan decomposition, connecting map}
\hypersetup{pdftitle={A strong Lefschetz family among equigenerated monomial almost complete intersections},pdfauthor={Ke-kuk Chow}}

\begin{document}
\begin{abstract}
Let $k$ be a field of characteristic zero and let $a\ge4$. We prove that
the algebra
\[
 k[x,y,z]/(x^a,y^a,z^a,x^{a-2}yz)
\]
has the strong Lefschetz property.
This establishes the second positive infinite family in the
Gasanova--Lundqvist--Nicklasson conjecture on equigenerated monomial almost
complete intersections of codimension three. An actual graded module
extension reduces all possible rank conflicts to two multiplication maps.
Explicit primitive vectors show that both connecting maps are nonzero.
We also determine that the sums of the endpoints of the graded Jordan
intervals take two consecutive values. The resulting almost-centered
structure, together with the standard tensor mechanism, proves the same
property after adjoining arbitrary truncated polynomial variables.
In particular, it proves the corresponding expectation in any number of
variables. The exclusion direction of the full classification conjecture
is outside the scope of the result.
\end{abstract}
\maketitle

\section{Introduction and precise scope}

For a standard graded Artinian algebra $B$ over a field $k$, a linear form
$L\in B_1$ is a \emph{strong Lefschetz element} if
\[
 L^{j-i}:B_i\longrightarrow B_j
\]
has maximal rank for every $0\le i<j$. The algebra has the strong
Lefschetz property, abbreviated SLP, if such an element exists.
The SLP of monomial complete intersections is classical
\cite{Stanley1980}. A natural next class consists of monomial almost
complete intersections, obtained by adding one monomial generator to the
ideal of pure powers.

Gasanova, Lundqvist, and Nicklasson \cite{GLN2022} developed a basis
splitting method for Lefschetz properties. In their Theorem~5 they proved
that
\[
 k[x_1,\ldots,x_n]/(x_1^a,\ldots,x_n^a,x_1^{a-1}x_2)
\]
has the SLP in characteristic zero. Immediately before Conjecture~1 in
Section~5.1 of version~4 of that preprint, they state the expectation that replacing
the additional generator by $x_1^{a-2}x_2x_3$, for $a\ge4$, also gives the
SLP. In codimension three this is the second positive infinite family
listed in that conjecture.

\begin{theorem}\label{thm:main}
Let $k$ be a field of characteristic zero and let $a\ge4$. In
\[
 A_a=k[x,y,z]/(x^a,y^a,z^a,x^{a-2}yz),
\]
the element $L=x+y+z$ is strong Lefschetz. Moreover, every graded Jordan
interval for multiplication by $L$ has endpoint sum either $3a-5$ or
$3a-4$.
\end{theorem}

\begin{corollary}\label{cor:extra}
Let $a\ge4$, $n\ge3$, and $d_4,\ldots,d_n\ge1$. Over a field of
characteristic zero, the sum of the variables is a strong Lefschetz
element of
\[
 \frac{k[x_1,\ldots,x_n]}
 {(x_1^a,x_2^a,x_3^a,x_4^{d_4},\ldots,x_n^{d_n},x_1^{a-2}x_2x_3)}.
\]
In particular, $k[x_1,\ldots,x_n]/(x_i^a,x_1^{a-2}x_2x_3)$ has the SLP.
\end{corollary}

Theorem~\ref{thm:main} establishes one specified positive family; it does
not establish the assertion that the families listed in
\cite[Conjecture~1]{GLN2022} exhaust all equigenerated examples with the
SLP. No assertion about positive characteristic is made here.

The general tools used in the proof are established ones: graded Jordan
decomposition, the lowering operator for monomial complete intersections,
and the connecting map of a module extension. The specific step is the
uniform nonvanishing calculation for two connecting maps. This calculation
is needed because splitting the Hilbert function does not split the
operator. The tensor argument in Corollary~\ref{cor:extra} is the usual
almost-centered mechanism, associated with Lindsey's work and recalled in
\cite[Theorem~2.5]{CJK2026}; we include a proof for the particular endpoint
condition used here.

Recent classification results in \cite{CJK2026} cover additional monomials
supported on two variables, or monomial almost complete intersections with
symmetric Hilbert series. The generator in Theorem~\ref{thm:main} has
support three, and its Hilbert series is not symmetric; see
Remark~\ref{rem:nonsymmetric}. Thus those hypotheses do not directly
give the result proved here.

\section{Graded strings and complete intersections}

Let $V$ be a finite graded vector space and let $U$ have degree one.
Regarding $V$ as a graded $k[u]$-module gives its graded Jordan
decomposition. A cyclic summand with a generator in degree $b$ and last
nonzero iterate in degree $e$ will be denoted by $[b,e]$. The integer
$b+e$ is its \emph{endpoint sum}. Intervals are counted with multiplicity.

\begin{lemma}\label{lem:nested}
All powers of $U$ have maximal rank if and only if its graded Jordan
intervals are totally ordered by containment. A nested interval family is
uniquely determined by its Hilbert function. In particular, a family whose
endpoint sums belong to $\{E,E+1\}$ is nested.
\end{lemma}

\begin{proof}
On one interval, a map from degree $i$ to degree $j>i$ has a kernel
contribution precisely when the interval contains $i$ but not $j$, and a
cokernel contribution precisely when it contains $j$ but not $i$.
Both kinds cannot occur in a nested family. Conversely, incomparable
intervals $[b,e]$ and $[c,f]$, with $b<c$ and $e<f$, give both kinds for
the map from degree $b$ to degree $f$. This proves the rank criterion.

If $h_j=\dim V_j$, the $r$th interval, ordered from the outside to the
inside, is the set $\{j:h_j\ge r\}$. This proves uniqueness. Finally,
if $b<c$ and both endpoint sums lie in $\{E,E+1\}$, then
\[
 e\ge E-b\ge E+1-c\ge f.
\]
Thus the intervals are comparable; intervals with equal left endpoint
are comparable as well.
\end{proof}

For positive integers $r_1,\ldots,r_s$, write
\[
 \CI(r_1,\ldots,r_s)
 =k[z_1,\ldots,z_s]/(z_1^{r_1},\ldots,z_s^{r_s}).
\]

\begin{lemma}[Classical complete-intersection strings]\label{lem:ci}
In characteristic zero, the graded Jordan intervals for multiplication by
$U=z_1+\cdots+z_s$ on $\CI(r_1,\ldots,r_s)$ all have endpoint sum
$E=\sum_i(r_i-1)$. On a shifted copy with shift $b$, the endpoint sum
is $E+2b$.
\end{lemma}

\begin{proof}
We recall the lowering-operator proof. First work over $\R$, put
$e_i=r_i-1$, and make the monomial basis orthogonal with
\[
 \|z^{\mathbf v}\|^2
 =\prod_i\frac{v_i!e_i!}{(e_i-v_i)!}.
\]
Define
\[
 V(z^{\mathbf v})
 =\sum_{i:v_i>0}v_i(e_i-v_i+1)z^{\mathbf v-\mathbf e_i}.
\]
Then $V=U^*$ and $VU-UV=(E-2j)\operatorname{id}$ on degree $j$.
If $Vf=0$ and $\deg f=b$, induction gives
\begin{equation}\label{eq:primitive}
 VU^qf=q(E-2b-q+1)U^{q-1}f.
\end{equation}
Consequently
\begin{equation}\label{eq:norm}
 \|U^qf\|^2=q(E-2b-q+1)\|U^{q-1}f\|^2.
\end{equation}
A nonzero such $f$ must have $w=E-2b\ge0$: otherwise the case $q=1$
contradicts positivity. Its iterates through $q=w$ are nonzero, and
$U^{w+1}f=0$. Thus it generates $[b,E-b]$.

Choose a vector in the least nonzero degree. It is killed by $V$.
The span of its string is preserved by $U$ and $V$, so its orthogonal
complement is also preserved by both operators and is graded. Induction
gives a decomposition entirely into strings $[b,E-b]$.

All multiplication matrices have integer entries. Their ranks over
$\R$ and $\Q$ agree, and every characteristic-zero field has the same
ranks. The ranks of all degree-to-degree maps determine the graded Jordan
multiplicities, so the endpoint assertion transfers to every such field.
The assertion about shifts is immediate.
\end{proof}

We will also use the following direct consequence of
\eqref{eq:primitive}--\eqref{eq:norm}. In a monomial complete intersection
of socle degree $E$, a homogeneous vector $f$ of degree $i$, killed by the
lowering operator, satisfies
\begin{equation}\label{eq:primitive-kernel}
 U^{E-2i+1}f=0\qquad\text{if }E-2i\ge0.
\end{equation}
The vectors used below have rational coefficients, so this identity also
holds over every characteristic-zero field.

\section{An actual extension and the two possible conflicts}

For the remainder of the three-variable proof, abbreviate
\[
 A=A_a,\qquad m=a-2,\qquad D=3a-5,\qquad t=y+z.
\]
Set
\[
 T=(x^m)A,\qquad
 C=A/T=\CI(a-2,a,a).
\]
These give an exact sequence of graded $k[L]$-modules
\begin{equation}\label{eq:extension}
 0\longrightarrow T\longrightarrow A\longrightarrow C\longrightarrow0.
\end{equation}
No splitting of this extension is assumed.

\begin{lemma}\label{lem:T}
The graded Jordan intervals of $T$ for multiplication by $L$ are
\begin{equation}\label{eq:Tstrings}
 [a-2,2a-2],\quad [a-1,2a-3],\quad
 [a-1,2a-2],\quad [a,2a-3].
\end{equation}
The first two have endpoint sum $D+1$ and the last two have endpoint
sum $D+2$.
\end{lemma}

\begin{proof}
The annihilator of $x^m$ gives the actual module description
\[
 T\cong x^m\bigl(k[x]/(x^2)\otimes_k
             k[y,z]/(yz,y^a,z^a)\bigr),
\]
where $x^m$ denotes a degree shift and $L$ acts as the sum of the
factor operators. In the second factor there is an invariant direct sum
\begin{equation}\label{eq:axes}
 \frac{k[y,z]}{(yz,y^a,z^a)}
 \cong k[t]/(t^a)\ \oplus\ yk[t]/(t^{a-1}).
\end{equation}
Indeed, the displayed summands have bases
$1,t,\ldots,t^{a-1}$ and $y,yt,\ldots,yt^{a-2}$.
For $r\ge1$, $t^r=y^r+z^r$ and $yt^{r-1}=y^r$;
these vectors form a basis in every degree.
For $r\ge2$, the two strings of $\CI(2,r)$ are $[0,r]$ and
$[1,r-1]$: their endpoint sums are $r$ and their Hilbert function is
$(1,2,\ldots,2,1)$. Tensoring the two strings in \eqref{eq:axes}
with $k[x]/(x^2)$, and applying the shifts $m$ and $m+1$,
respectively, now yields \eqref{eq:Tstrings}.
\end{proof}

Every string of $C$ has endpoint sum $D$. For $i<j$, if $i+j\le D$,
the maps $L^{j-i}$ on both $T$ and $C$ are injective, and hence so is the
map on $A$. If $i+j\ge D+2$, both maps are surjective, and so is the
map on $A$. Only the line $i+j=D+1$ remains. On that line the map on
$C$ is surjective. The strings of $T$ with endpoint sum $D+1$ do not
give a cokernel. For a string $[b,e]$ of endpoint sum $D+2$, a cokernel
on this line requires $i<b$ and $j\le e$; the identity
$i+j=b+e-1$ forces $i=b-1$ and $j=e$.
There are therefore exactly two maps requiring further work:
\begin{equation}\label{eq:two}
 (i,j,k=j-i)=(a-2,2a-2,a),\qquad
             (a-1,2a-3,a-2).
\end{equation}

In each case the $T$-cokernel has dimension one. The $C$-kernel has
dimension $a-2$. To see the latter, symmetry gives
$h_C(j)=h_C(i-1)$; direct counting in degrees $a-2$ and $a-1$ gives
\[
 h_C(a-2)-h_C(a-3)=h_C(a-1)-h_C(a-2)=a-2.
\]
Applying the snake lemma to \eqref{eq:extension} gives
\begin{equation}\label{eq:delta}
 \ker(L^k:C_i\to C_j)\xrightarrow{\ \delta\ }
 \coker(L^k:T_i\to T_j)\longrightarrow
 \coker(L^k:A_i\to A_j)\longrightarrow0,
\end{equation}
where
\[
 \delta(f)=[L^k\widetilde f]
\]
for any lift $\widetilde f\in A_i$ of $f$. Thus a nonzero value of
$\delta$ proves surjectivity on $A$ at each of the two maps.

\section{Primitive vectors and the connecting maps}

For $0\le r\le a-1$, define a homogeneous binary polynomial by
\begin{equation}\label{eq:Pr}
 P_r(y,z)=\sum_{\ell=0}^r c_\ell y^{r-\ell}z^\ell,
 \qquad c_0=1,\qquad
 c_{\ell+1}=-c_\ell
 \frac{(r-\ell)(a-r+\ell)}{(\ell+1)(a-1-\ell)}.
\end{equation}
The denominators are nonzero in characteristic zero.

\begin{lemma}\label{lem:Pr}
In $\CI(a,a)$, $(V_y+V_z)P_r=0$, and the coefficients of $y^r$ and
$z^r$ are $1$ and $(-1)^r$, respectively. In $C$, put
$\mu=a-3$ and $\lambda=2a-2-2r$. Both $P_r$ and
\begin{equation}\label{eq:coupled}
 \lambda xP_r-\mu tP_r
\end{equation}
are killed by $V_x+V_y+V_z$ whenever the latter expression is used below.
\end{lemma}

\begin{proof}
The coefficient of $y^{r-\ell-1}z^\ell$ in
$(V_y+V_z)P_r$ is
\[
 c_\ell(r-\ell)(a-r+\ell)
 +c_{\ell+1}(\ell+1)(a-1-\ell),
\]
which is zero by \eqref{eq:Pr}. Taking the product of the successive
ratios also gives $c_r=(-1)^r$. The $x$-factor of $C$ has top exponent
$a-3$, so $V_xx=\mu$. Equation~\eqref{eq:primitive} in the binary
factor gives $(V_y+V_z)tP_r=\lambda P_r$.
These two contributions cancel in \eqref{eq:coupled}.
\end{proof}

\subsection{\texorpdfstring{The map from degree $a-2$ to degree $2a-2$}{The first connecting map}}

Here $k=a$. The source $T_{a-2}$ is generated by $x^m$, and the target
has basis
\[
 u=x^{a-1}y^{a-1},\qquad v=x^{a-1}z^{a-1}.
\]
The image of $x^m$ is $a(u+v)$. A functional detecting its
one-dimensional cokernel is therefore
\[
 \ell_1(u)=1,\qquad \ell_1(v)=-1.
\]

If $a$ is odd, choose $f=P_{a-2}$. It is primitive in $C$ of degree
$a-2$, so \eqref{eq:primitive-kernel} gives $L^af=0$ in $C$.
Its pure-axis terms are $y^{a-2}-z^{a-2}$.
Only these terms can contribute to $u$ or $v$ after multiplication.
Their coefficients in $L^af\in T$ are $a$ and $-a$, whence
\begin{equation}\label{eq:delta1odd}
 \ell_1(L^af)=2a\ne0.
\end{equation}

If $a$ is even, choose
\[
 f=4xP_{a-3}-(a-3)tP_{a-3}.
\]
Lemma~\ref{lem:Pr} applies with $\lambda=4$; the same kernel identity
holds. The terms on the $y$-axis are
$4xy^{a-3}-(a-3)y^{a-2}$, and the terms on the $z$-axis are their
negatives with $y$ replaced by $z$. Hence the coefficient of $u$ in
$L^af$ is
\[
 4\binom a2-(a-3)a=a(a+1),
\]
and that of $v$ is its negative. Thus
\begin{equation}\label{eq:delta1even}
 \ell_1(L^af)=2a(a+1)\ne0.
\end{equation}
Equations~\eqref{eq:delta1odd} and \eqref{eq:delta1even} show that
the first connecting map in \eqref{eq:two} is surjective.

\subsection{\texorpdfstring{The map from degree $a-1$ to degree $2a-3$}{The second connecting map}}

Now put $k=a-2$. In $T_{2a-3}$ use the ordered basis
\begin{align*}
 u&=x^{a-2}y^{a-1},&v&=x^{a-2}z^{a-1},\\
 p&=x^{a-1}y^{a-2},&q&=x^{a-1}z^{a-2}.
\end{align*}
The source basis $x^{a-1},x^{a-2}y,x^{a-2}z$ maps to
\[
 p+q,\qquad u+kp,\qquad v+kq,
\]
respectively. These vectors are independent, and a functional on the
cokernel is
\begin{equation}\label{eq:ell2}
 \ell_2(u,v,p,q)=(-k,k,1,-1).
\end{equation}

If $a$ is even, choose $f=P_{a-1}$. It is primitive of degree $a-1$
in $C$, and \eqref{eq:primitive-kernel} gives $L^{a-2}f=0$ in $C$.
Its pure-axis terms are $y^{a-1}-z^{a-1}$. The coefficients of
$(u,v,p,q)$ in $L^{a-2}f$ are $(1,-1,0,0)$, so
\begin{equation}\label{eq:delta2even}
 \ell_2(L^{a-2}f)=-2(a-2)\ne0.
\end{equation}

If $a$ is odd, then $a\ge5$ and we choose
\[
 f=2xP_{a-2}-(a-3)tP_{a-2}.
\]
Lemma~\ref{lem:Pr} now applies with $\lambda=2$.
The degree is $a-1$, so again $L^{a-2}f=0$ in $C$.
The $y$-axis terms are
$2xy^{a-2}-(a-3)y^{a-1}$; the $z$-axis terms are their negatives.
The $u$-coefficient of $L^kf$ is $2k-(a-3)=a-1$, and the
$p$-coefficient is $2$. Consequently
\begin{equation}\label{eq:delta2odd}
 \ell_2(L^kf)=2\bigl(-(a-2)(a-1)+2\bigr)=-2a(a-3)\ne0.
\end{equation}
Thus the second connecting map is surjective as well.

\begin{proof}[Proof of the SLP assertion in Theorem~\ref{thm:main}]
Every map except the two in \eqref{eq:two} has already been proved
injective or surjective by the argument following Lemma~\ref{lem:T}.
The four calculations above prove that the connecting maps for the
remaining two cases are nonzero and hence surjective. Exactness of
\eqref{eq:delta} proves surjectivity on $A$ in those cases. All powers
therefore have maximal rank.
\end{proof}

\section{Endpoint control and additional variables}

To prove the remaining assertion of Theorem~\ref{thm:main}, we determine
enough of the actual graded Jordan intervals. This step uses the SLP
just proved; a Hilbert-function calculation on its own would not suffice.

\begin{proposition}\label{prop:endpoints}
The graded Jordan intervals of $A$ have endpoint sums in $\{D,D+1\}$.
\end{proposition}

\begin{proof}
Each interval of $C$ has the form $[p,D-p]$.
At each of $p=a-2$ and $p=a-1$ its multiplicity is $a-2$, by the
Hilbert-function differences computed in Section~3.
For each of these two values of $p$, take one such interval and the
corresponding interval $[p+1,D-p+1]$ of $T$ with endpoint sum $D+2$.
There is an identity of interval indicator functions
\begin{equation}\label{eq:rearrange}
 [p,D-p]+[p+1,D-p+1]
 =[p,D-p+1]+[p+1,D-p].
\end{equation}
The intervals on the right are nonempty for $a\ge4$; the smallest
intersection occurs at $a=4,p=3$ and is $[4,4]$.

After the two replacements, all untouched $C$ intervals have endpoint
sum $D$, and all other intervals have endpoint sum $D+1$.
They form a nested family by Lemma~\ref{lem:nested}, and their Hilbert
function is that of $A$, by \eqref{eq:extension} and
\eqref{eq:rearrange}. The actual intervals of $A$ are also nested,
because the SLP has been proved. Uniqueness in Lemma~\ref{lem:nested}
therefore identifies the actual interval family with the constructed
family and proves the assertion.
\end{proof}

\begin{proof}[Proof of Corollary~\ref{cor:extra}]
Put
\[
 B=k[x_4,\ldots,x_n]/(x_4^{d_4},\ldots,x_n^{d_n}),
 \qquad E=\sum_{i=4}^n(d_i-1).
\]
The algebra in the corollary is $A\otimes_k B$. Tensor each actual
string $[b,e]$ of $A$ with $B$. As a module for the sum of the
operators this is a shifted monomial complete intersection; hence
Lemma~\ref{lem:ci} decomposes it into strings with endpoint sum
$b+e+E$. Proposition~\ref{prop:endpoints} says that all final endpoint
sums belong to $\{D+E,D+E+1\}$. They are therefore nested by
Lemma~\ref{lem:nested}, which proves the SLP.
\end{proof}

\begin{remark}\label{rem:nonsymmetric}
The socle degree of $A$ is $D$. Its Hilbert function is not symmetric:
$C$ is symmetric about $D/2$, whereas
\[
 \dim T_{a-2}=1,\qquad \dim T_{2a-3}=4.
\]
Since $D-(a-2)=2a-3$, this gives
\[
 h_A(a-2)-h_A(2a-3)=-3.
\]
Thus the symmetric-Hilbert-series results of
\cite[Proposition~4.2 and Theorem~4.3]{CJK2026} do not directly apply.
Endpoint control, rather than symmetry, is what allows the tensor
corollary here.
\end{remark}

\section{Reproducibility and remaining scope}

The argument is independent of computation. Ancillary files provide
finite checks in the original monomial coordinates. For $a=4,\ldots,20$,
the code constructs the primitive vector for each of the two critical
maps, checks its vanishing in $C$, and checks its nonzero connecting
value in $A$, giving $34$ certificates. It also constructs every positive
power map for
\[
 (n,a)=(3,4),(3,5),(3,6),(4,4),(4,5),
\]
and computes their exact ranks over $\Q$ using rational Gaussian
elimination. These are $334$ degree-to-degree maps. Graded Jordan
intervals are recovered from these ranks by inclusion-exclusion.
The checks are supplied to facilitate reproduction and attempts to
falsify the formulas, and are not used to infer the universal theorem.

The result proves the second positive infinite family of
\cite[Conjecture~1]{GLN2022}, in the numbering of
arXiv:2006.14453v4, and the preceding expectation in arbitrary dimension.
It leaves the exclusion of other exponent patterns in that full
classification untouched. It also does not give a positive-characteristic
classification. The use of characteristic zero is explicit in both the
complete-intersection strings and the nonvanishing calculations.

\section*{AI assistance disclosure}
This manuscript was developed with substantial assistance from OpenAI
Codex for mathematical derivations, drafting, and computational checks.
The AI system is not an author. This disclosure does not constitute
external mathematical review or certification.

\bibliographystyle{amsplain}
\bibliography{references}
\end{document}
